\section{Companion lab: validate UTF-8 across chunks}
\label{revision-lab:utf8}

The packet parser can leave an incomplete request untouched because the caller
still owns the whole prefix. A stream interface may instead accept that prefix
now and receive later bytes in another call. Its state must remember what the
accepted bytes permit next. Resetting at each call would make correctness depend
on where a file reader happened to divide the input.

This companion develops original project \texttt{015-utf8-scanner} through the
native \texttt{015-utf8-streams} unit. Open its assignment and begin with the
complete \texttt{worked/range.asm} experiment. The second complete source,
\texttt{worked/lower-bound.asm}, changes one sample byte so that you can compare
rejection just below a bound with acceptance at that bound.

Implement \texttt{scan\_utf8} in \texttt{independent/stream.asm}. Its allocations,
chunk caller and \texttt{chunk\_ready} observation marker are complete. Replace the
single deliberate \texttt{break} body while retaining that public interface.
Choose your own temporary registers, private labels and equivalent algorithm.
The three inline assessments check direct chunk transitions, interleaved state
objects and the supplied caller. The archived construction regions are not
required by this assignment.

The prerequisite is Chapter 12's bounded byte access and function reasoning.
The new work is a persistent state machine: define which byte can follow a prefix,
carry that decision across calls, count completed values, and retain the first
failure. This scanner validates and counts UTF-8. It does not produce decoded
text or replacement characters.

\subsection*{Distinguish bytes, scalar values and displayed text}

A Unicode code point is a numbered position, conventionally written with
\texttt{U+} and hexadecimal digits. A Unicode \emph{scalar value} is a code point
outside the surrogate range U+D800 through U+DFFF. UTF-8 represents a scalar with
one through four bytes. The restrictions in Table~\ref{tab:revision-utf8-leads}
exclude surrogate encodings, values beyond U+10FFFF and encodings longer than
needed.\footnote{The Unicode Consortium, \emph{The Unicode Standard}, version
17.0.0, Section 3.9.3 and Table 3-7:
\url{https://www.unicode.org/versions/Unicode17.0.0/core-spec/chapter-3/}.}

A complete scalar is the unit this lab counts. It is not necessarily one visible
character. For example, U+0065 followed by U+0301 represents an \texttt{e} and a
combining acute accent: two scalars that can form one grapheme cluster in displayed
text. U+00E9 is one scalar with a precomposed accented letter. This scanner counts
them as two and one respectively; it does not normalize either representation or
segment grapheme clusters.\footnote{For grapheme clusters and text boundaries,
see Unicode Standard Annex 29, revision 47, Section 3:
\url{https://www.unicode.org/reports/tr29/tr29-47.html}.}

The supplied stream has eleven bytes, written here in hexadecimal:
\[
\underbrace{\mathtt{41}}_{\mathrm{U+0041}}\quad
\underbrace{\mathtt{E2\ 82\ AC}}_{\mathrm{U+20AC}}\quad
\underbrace{\mathtt{F0\ 9F\ 99\ 82}}_{\mathrm{U+1F642}}\quad
\underbrace{\mathtt{00}}_{\mathrm{U+0000}}\quad
\underbrace{\mathtt{C3\ A9}}_{\mathrm{U+00E9}}.
\]
It contains five complete scalars. Byte zero is a valid one-byte encoding; it
neither ends this supplied stream nor suppresses the following scalar. U+FEFF,
unassigned scalar values and noncharacters also count when their encodings are
well formed. Assignment and display policies are outside this validator's contract.
The program exits silently, so inspect its state rather than waiting for printed
characters.

\subsection*{Derive the next-byte rule from an accepted prefix}

A one-byte encoding has high bit zero. A multibyte lead specifies its width;
following bytes, called \emph{continuation bytes}, have bit pattern
\texttt{10xxxxxx}. Their ordinary hexadecimal interval is \texttt{80} through
\texttt{BF}. A continuation cannot begin a new scalar. Bytes \texttt{C0},
\texttt{C1} and \texttt{F5} through \texttt{FF} cannot begin a valid encoding
either. The first continuation after four particular leads has a tighter bound.

\begin{table}[htbp]
  \centering
  \small
  \begin{tabular}{llll}
    \toprule
    Lead bytes (hex) & Width & First continuation (hex) & Later continuations (hex) \\
    \midrule
    00--7F           & 1     & none                     & none                      \\
    C2--DF           & 2     & 80--BF                   & none                      \\
    E0               & 3     & A0--BF                   & 80--BF                    \\
    E1--EC, EE--EF   & 3     & 80--BF                   & 80--BF                    \\
    ED               & 3     & 80--9F                   & 80--BF                    \\
    F0               & 4     & 90--BF                   & 80--BF                    \\
    F1--F3           & 4     & 80--BF                   & 80--BF                    \\
    F4               & 4     & 80--8F                   & 80--BF                    \\
    \bottomrule
  \end{tabular}
  \caption{UTF-8 lead classes and the ranges used by this scanner. Width includes the lead. Only the first continuation can have a restricted range; later continuations use the ordinary range. Derived from Unicode 17.0.0 Table 3-7.}
  \label{tab:revision-utf8-leads}
\end{table}


These bounds do more than check the continuation tag. After \texttt{E0}, accepting
\texttt{80} would allow a value below the three-byte minimum. After \texttt{ED},
accepting \texttt{A0} would enter the surrogate range. The \texttt{F0} lower bound
excludes values that need only three bytes; the \texttt{F4} upper bound excludes
values above the Unicode limit. After a valid first continuation, reset the range
to \texttt{80} through \texttt{BF} for every remaining continuation. The tighter
rule does not persist for the whole sequence.\footnote{The byte grammar is also
specified in RFC 3629, Section 4:
\url{https://www.rfc-editor.org/rfc/rfc3629.html\#section-4}.}

You can derive a boundary numerically. A three-byte encoding carries four useful
bits from the lead and six from each continuation. For bytes $b_0,b_1,b_2$, its
value is
\[
 ((b_0\mathbin{\&}15)\ll12)
 \mathbin{|}((b_1\mathbin{\&}63)\ll6)
 \mathbin{|}(b_2\mathbin{\&}63).
\]
Thus \texttt{E0 A0 80} gives $32\cdot64=2048$, hexadecimal 0800, the first
three-byte scalar. \texttt{E0 9F 80} would give 1984, which has a shorter encoding.
The scanner need not assemble this numerical value: retaining the allowed range
and remaining count is enough to validate it. Use the formula to explain the
range before encoding comparisons in assembly.

Load input as unsigned bytes when comparing these ranges. An \texttt{lb} of
\texttt{E0} sign-extends a negative value; unsigned comparisons with small positive
bounds would then classify the wrong number. Advance the byte pointer with
full-width arithmetic even though this lab's byte counts fit within a word.
Neither a short chunk nor a small scalar implies a small address.

\subsection*{Isolate an inclusive byte comparison}

The complete range experiment loads byte \texttt{9F} and calls this helper with
lower bound \texttt{A0} and upper bound \texttt{BF}. The arguments are ordinary
unsigned byte values in registers, so decimal 159, 160 and 191 describe the same
comparison. The helper returns one for membership and zero otherwise. Its result
is separate from the scanner's status codes.

\noindent\begin{minipage}{\linewidth}
\MARSCodeLabel{Test a byte against inclusive unsigned bounds}
\begin{minted}{mars-idea}
within_range:
    sltu $t0, $a0, $a1
    bnezc $t0, outside_range
    nop
    sltu $t0, $a2, $a0
    xori $v0, $t0, 1
    jrc $ra
outside_range:
    move $v0, $zero
    jrc $ra
\end{minted}
\end{minipage}

The first \texttt{sltu} produces one if the byte is below the lower bound. The
compact branch then selects the zero-return path. If that test is zero, the second
comparison asks whether the upper bound is below the byte. That comparison is one
only for a byte that is too large; \texttt{xori} reverses its zero-or-one result.
Equality with either bound therefore remains accepted.

Use the first guided Debug sequence to observe \texttt{9F}: the lower-bound test
is one and the returned result is zero. In \texttt{worked/lower-bound.asm}, the
sample becomes \texttt{A0}. Its lower-bound test is zero and the returned result
is one. Both callers store the answer in \texttt{range\_result} and exit silently.
The public \texttt{worked/range.kts} script checks every byte value against seven
relevant intervals as well as the first complete caller.

A correct interval helper is only one part of a scanner. The accepted prefix
chooses the interval, the pending count determines whether a scalar is complete,
and the call's final flag decides what an unfinished prefix means at the end.
Those decisions need persistent state, which the helper deliberately does not
provide.

\subsection*{Give persistent state a precise meaning}

The function takes the selected state pointer in \texttt{\$a0}, current chunk
pointer in \texttt{\$a1}, byte count in \texttt{\$a2}, and final flag in
\texttt{\$a3}. The flag is zero or one. A chunk contains zero through 64 bytes
inside the supplied 96-byte arena and can begin at any byte offset that leaves
room for it. State is a separate, word-aligned 32-byte object, either
\texttt{state\_a} or \texttt{state\_b}. Calls receive a reachable state; you do
not need to accept arbitrary corrupted descriptors. For an active call, old total
plus supplied count is at most 4096.

\begin{table}[htbp]
\centering
\small
\begin{tabular}{rll}
\toprule
Byte offset & State word & Meaning at return\\
\midrule
0 & total & Accepted bytes over all calls\\
4 & count & Completed scalars over all calls\\
8 & need & Continuations still needed\\
12 & minimum & Lowest permitted next continuation\\
16 & maximum & Highest permitted next continuation\\
20 & start & Pending or failed lead's absolute offset; $-1$ if idle\\
24 & status & 0 active, 1 finished, $-1$ malformed, $-2$ truncated\\
28 & error & Detecting offset, or $-1$ before any failure\\
\bottomrule
\end{tabular}
\caption{The selected scanner state consists of eight native words. Total counts accepted bytes, so it also names the offset where the next byte would be accepted. A terminal rejection returns $-3$ without replacing the stored status.}
\label{tab:revision-utf8-state}
\end{table}


The empty state is $[0,0,0,128,191,-1,0,-1]$, in the table's word order and decimal
notation. An idle active state has need zero, ordinary continuation bounds and
start $-1$. A pending active state has need from one through three and retains
the absolute offset of its lead. Its bounds describe the \emph{next} byte, not
the first byte of the next call. Error remains $-1$ while active.

At each input-byte boundary, total is the number of accepted bytes over all calls;
count is the number of complete scalars among them. A valid pending prefix
contributes to total but not yet to count. The current call's accepted count is
separate. Return it in \texttt{\$v1}; total includes earlier calls too.

If need is zero, an ASCII byte completes a scalar immediately. A legal multibyte
lead records start at the old total, establishes need and the first-continuation
range, and is then accepted. If need is positive, test the current byte against
the stored range. An accepted continuation reduces need and restores ordinary
bounds. Reaching need zero completes one scalar and resets start to $-1$.
Only acceptance advances total, the call's accepted count and the input pointer.
An impossible byte advances none of them.

Keep this invariant at the loop's top. First determine whether the current call
has an unexamined byte, then load within its supplied interval. A zero-length
call performs no input load. The contract permits read-ahead within the current
chunk, but never across its end, even if the next chunk is mapped immediately
after it. A chunk ending with \texttt{E2} does not grant access to the following
\texttt{82}; retain the pending rule for the next call instead.

\subsection*{Trace one scalar across two chunk boundaries}

The caller supplies the eleven-byte example as four chunks, with the last one
marked final. Predict every row of Table~\ref{tab:revision-utf8-chunks} before
running your implementation. All offsets are zero based; state and return values
are decimal unless the byte column says otherwise.

\begin{table}[htbp]
  \centering
  \small
  \setlength{\tabcolsep}{5pt}
  \begin{tabular}{rlrrrrll}
    \toprule
    Index & Bytes (hex) & Total & Count & Need & Start & Next range & $(v_0,v_1)$ \\
    \midrule
    0     & 41 E2       & 2     & 1     & 2    & 1     & 128--191   & $(0,2)$     \\
    1     & 82 AC F0    & 5     & 2     & 3    & 4     & 144--191   & $(0,3)$     \\
    2     & 9F 99       & 7     & 2     & 1    & 4     & 128--191   & $(0,2)$     \\
    3     & 82 00 C3 A9 & 11    & 5     & 0    & $-1$  & 128--191   & $(1,4)$     \\
    \bottomrule
  \end{tabular}
  \caption{Completed calls for the supplied four-chunk stream. Only index three is final. Stored status equals $v_0$ in these active calls; error stays $-1$. The returned accepted count is local to each call, while total and count accumulate.}
  \label{tab:revision-utf8-chunks}
\end{table}


The first chunk accepts \texttt{A} and the \texttt{E2} lead. Total becomes two,
but only one scalar is complete. The second chunk finishes U+20AC and accepts
\texttt{F0}, so need becomes three and the next-byte minimum becomes 144,
hexadecimal \texttt{90}. The third chunk accepts two continuations. The first,
\texttt{9F}, satisfies that tighter bound; after it, minimum returns to 128.
Need is still one when the chunk ends. Its two bytes have not completed another
scalar.

The fourth chunk begins with \texttt{82}, which completes U+1F642. It then accepts
the zero byte and \texttt{C3 A9}. Count becomes five, need becomes zero and the
final flag permits successful completion. Splitting the same bytes at another
position must produce the same final total and scalar count, although intermediate
states and per-call accepted counts will differ.

Break at \texttt{chunk\_ready} in Debug. The caller has already stored
\texttt{chunk\_index}, \texttt{last\_status} and \texttt{last\_used}; their values
identify the completed call. The state words must agree with that same call.
For the native Debug activity, place the breakpoint on the \texttt{nop} beneath
\texttt{chunk\_ready} and set its condition to \texttt{\$s3 == 0}. After that
stop, change only the condition to \texttt{\$s3 == 1}, then resume. Repeat for
2 and 3. Keep the same execution session and source throughout: restarting would
return to the first chunk, while editing source would invalidate the observation. The revised caller uses
compact transfers and preserves its live values in saved registers. Your scanner
must preserve those registers and the stack pointer across all returns.

Figure~\ref{fig:revision-utf8-pending} compares the pending rules after chunks
one and two in actual native Memory observations. Read the boxed words as need
and minimum using the state layout above.

\begin{figure}[htbp]
\centering
\begin{tikzpicture}[alt={Actual Memory observations at chunk_ready after chunks one and two. The first state begins with words total five, count two, need three and minimum 144. The next begins with total seven, count two, need one and minimum 128. Both retain maximum 191, lead start four, active status zero and error minus one. Native boxes identify need and minimum.}]
\node[anchor=north west,inner sep=0,font=\small\bfseries] (atitle) at (0,0) {1. After chunk 1: a four-byte lead accepted};
\node[anchor=north west,inner sep=0] (a) at ([yshift=-2mm]atitle.south west) {\MARSRevisionCrop{revision/captures/utf8/utf8-four-byte-lead.png}{0.96\linewidth}{0.032812500}{0.314444444}{0.580468750}{0.608888889}};
\node[anchor=north west,inner sep=0,font=\small\bfseries] (btitle) at ([yshift=-4mm]a.south west) {2. After chunk 2: one continuation still needed};
\node[anchor=north west,inner sep=0] (b) at ([yshift=-2mm]btitle.south west) {\MARSRevisionCrop{revision/captures/utf8/utf8-pending-tail.png}{0.96\linewidth}{0.032812500}{0.314444444}{0.580468750}{0.608888889}};
\foreach \panel in {a,b} {
\begin{scope}[shift={(\panel.south west)},x={($(\panel.south east)-(\panel.south west)$)},y={($(\panel.north west)-(\panel.south west)$)}]
\draw[black,line width=.7pt] (.592,.319) rectangle (.988,.623);
\end{scope}}
\end{tikzpicture}
\caption{Actual light-theme states for the supplied stream. Boxes mark need and next minimum: three/144 becomes one/128, while count stays two. The first allowed continuation resets the tighter minimum. Each four-byte group is one native little-endian word.}
\label{fig:revision-utf8-pending}
\end{figure}


\subsection*{Separate malformed input, truncation and terminal reuse}

A nonfinal chunk that ends in a still-possible prefix returns zero. A final chunk
with no pending prefix returns one; a final chunk with a pending prefix returns
$-2$, meaning truncated. The latter is a promise that no more bytes will arrive,
not a new byte to validate. Even an empty final call can therefore finish or
truncate a stream.

An impossible byte returns $-1$ immediately, whether or not that call is final.
Store its absolute offset in error, retaining the lead offset in start. For an
invalid byte encountered while idle, that byte itself supplies start. Accepted
progress preceding the error stays committed. This differs from the packet
parser's no-output-on-failure contract: the stream's state records both its
accepted prefix and its terminal diagnostic.

Compare final input \texttt{41 E0} with final input \texttt{41 E0 80}. Both
accept two bytes and complete one scalar. The first ends with a possible prefix:
status $-2$, start one, error two. In the second, byte \texttt{80} at offset two
is already impossible after \texttt{E0}: status $-1$, start one, error two. Do not
call it truncated merely because the chunk was final. The offending byte is not
included in \texttt{\$v1} or total. A final flag never replaces an earlier byte
error. Figure~\ref{fig:revision-utf8-errors} compares those two completed calls.

\begin{figure}[htbp]
\centering
\begin{tikzpicture}[alt={Actual Memory observations for final chunks 41 E0 80 and 41 E0. Both states have total two, scalar count one, need two, minimum 160, maximum 191, start one and detecting offset two. The malformed state's status word is FFFFFFFF, minus one; the truncated state's is FFFFFFFE, minus two. Native boxes mark status and detecting offset.}]
\node[anchor=north west,inner sep=0,font=\small\bfseries] (atitle) at (0,0) {1. Final bytes 41 E0 80: malformed};
\node[anchor=north west,inner sep=0] (a) at ([yshift=-2mm]atitle.south west) {\MARSRevisionCrop{revision/captures/utf8/utf8-malformed-final.png}{0.96\linewidth}{0.032812500}{0.314444444}{0.580468750}{0.608888889}};
\node[anchor=north west,inner sep=0,font=\small\bfseries] (btitle) at ([yshift=-4mm]a.south west) {2. Final bytes 41 E0: truncated};
\node[anchor=north west,inner sep=0] (b) at ([yshift=-2mm]btitle.south west) {\MARSRevisionCrop{revision/captures/utf8/utf8-truncated-final.png}{0.96\linewidth}{0.032812500}{0.314444444}{0.580468750}{0.608888889}};
\foreach \panel in {a,b} {
\begin{scope}[shift={(\panel.south west)},x={($(\panel.south east)-(\panel.south west)$)},y={($(\panel.north west)-(\panel.south west)$)}]
\draw[black,line width=.7pt] (.592,.014) rectangle (.988,.319);
\end{scope}}
\end{tikzpicture}
\caption{The same accepted prefix has different terminal causes. Native boxes mark status and detecting offset. In the first picture, \texttt{FF FF FF FF} stores $-1$; in the second, \texttt{FE FF FF FF} stores $-2$. Both detect at offset two, but only the first has an impossible byte there. These are completed calls at \texttt{chunk\_ready}.}
\label{fig:revision-utf8-errors}
\end{figure}


A successful final call, a malformed call and a truncated final call all make
the state terminal. On any later call, return $(-3,0)$ without reading input or
writing caller data. Preserve the stored status and error, so an attempted reuse
does not erase the original result. Check status before entering the active scan.
A zero-length later call has the same rejection rule. To begin an unrelated stream,
the caller must deliberately initialize a fresh empty state.

While active, stores may change only the selected state. Preserve the input,
guards, other state, caller metadata and full saved registers. A correctly restored
private stack frame is permitted. Active store order and temporary registers are
not prescribed. A terminal rejection permits no caller-data stores, including a
write that restores the same value. Keep terminal rejection outside any common
active publication path that would rewrite state unconditionally.

\subsection*{Build tests around the first changed decision}

Begin with ASCII and an empty nonfinal call, then one multibyte scalar at each
width. Divide each valid encoding after every possible byte and insert empty
nonfinal chunks. Next test the four restricted leads on both sides of their first-
continuation bounds. Include \texttt{E0 A0 80} and \texttt{F0 90 80 80} to expose
a routine that forgets to restore ordinary bounds after the first continuation.

Test \texttt{ED A0} and \texttt{F4 90} before supplying any further byte: both
are already malformed. Test lone \texttt{80}, \texttt{C0} and \texttt{F5} while
idle. Repeat a pending prefix once with final zero and once with final one, then
try another call after each terminal outcome. Keep the original failure status
and offsets in view during that last experiment.

Use two state objects for interleaved independent streams. Feed \texttt{E2} to
one and \texttt{C3} to the other; then finish them with \texttt{82 AC} and
\texttt{A9} respectively. Their saved pending rules belong to those state objects,
not to private global variables or a register value left over from the last call.
Choose a fresh input offset for the second pair of calls. The supplied checks also
change scratch registers between calls, so register residue cannot serve as
persistent state.

Use \emph{Check this program} in the native lesson for the 426 original stream
cases, twelve additional interleaved-state cases and fifteen caller variants.
The interleaving assessment alternates both state objects within one case, checks
their complete state after each call, changes temporary registers and revisits
terminal outcomes. All three assessments use the same continuing source.

The two original public check domains are also available through the CLI. From
the downloaded unit's directory, run these commands on your completed source:

\noindent\begin{minipage}{\linewidth}
\begin{verbatim}
java -jar /path/to/cli.jar grade \
  independent/stream.kts independent/stream.asm
java -jar /path/to/cli.jar grade \
  independent/caller.kts independent/stream.asm
\end{verbatim}
\end{minipage}

The additional interleaving cases run through the native assessment. Finish with
an ordinary Run of the four-chunk caller; its successful silent exit is separate
from the state and access checks above.


When a check fails, retain the stream prefix, selected state, current chunk and
final flag. Find the first byte whose permitted range or accepted count diverges.
Then compare the saved state before and after that call. A correct final scalar
count cannot excuse an out-of-chunk load, an overwritten earlier error or lost
pending state. Keep one boundary trace and one malformed-versus-truncated trace
with your passing results, explaining why changing chunk boundaries does not
change the validity of the complete stream.
